% ===========================================================================
\section{Probability laws and surprise}\label{sec:surprise}
\warmth{0}\quad\whisper{The probability law is the object; surprise is a logarithm of that law.}

\begin{strand}
Before measuring information, keep the random variable separate from its law.
For a discrete random variable, the probability mass function records the chance
of every possible value. The same law may also be viewed as a probability measure
or through its cumulative distribution function.
\end{strand}

\block{Three views of the law}

\begin{definition}[probability mass function]\label{def:pmf}
Let $X$ take values in a discrete alphabet $\X$. Its probability mass function is
\[
  p:\X\longrightarrow[0,1],
  \qquad p(x)=\PP(X=x).
\]
We also write $p_X$ when the random variable must be made explicit.
\end{definition}

\begin{definition}[distribution and CDF]
The distribution of a random variable $X$ is denoted $P_X$. For any set $A$ in the value space, define
\[
  P_X(A)=\PP(X\in A).
\]
If $X$ is a real-valued random variable, its cumulative distribution function is defined as
\[
  F_X(y)=\PP(X\le y),\qquad y\in\mathbb{R}.
\]
\end{definition}

\begin{yourturn}
Complete the dictionary:
\[
  p_X(x)=\answerin[2.8cm]{\PP(X=x)},\qquad
  P_X((a,b))=\answerin[3.4cm]{\PP(a<X<b)}.
\]
\recall{$p_X(x)=\PP(X=x)$ and $P_X((a,b))=\PP(a<X<b)$.}
\end{yourturn}

\block{Why surprise is logarithmic}

\begin{definition}[surprise of an event]\label{def:surprise}
We seek a function $s(A)$ that depends continuously on $\PP(A)$, decreases as
$\PP(A)$ increases, and satisfies
\[
  s(A\cap B)=s(A)+s(B)
\]
whenever $A$ and $B$ are independent. Up to the choice of logarithm base, these
requirements give
\[
  s(A)=-\log\PP(A).
\]
\end{definition}

\begin{yourturn}
If independent events have probabilities $\PP(A)=1/2$ and $\PP(B)=1/4$, then
\[
  s(A)=\answerin[1.1cm]{1},\qquad
  s(B)=\answerin[1.1cm]{2},\qquad
  s(A\cap B)=\answerin[1.1cm]{3}\ \text{bits}.
\]
\end{yourturn}
\recall{$1$, $2$, and $3$ bits. Independence gives $\PP(A\cap B)=\PP(A)\PP(B)=1/8$.}
